\chapter{Couche 7 — Le contrôleur et les nœuds}

Le contrôleur doit créer et détruire des conteneurs sur des machines distantes.
C'est le lien le plus sensible de l'architecture : il porte, de bout en bout,
un pouvoir équivalent à \cmd{root}.

% -----------------------------------------------------------------------------
\mesure{SSH plutôt que l'API Docker}{mesure 24 · \fichier{workers.py}}

\menace
La façon évidente de piloter un démon Docker à distance est d'exposer son API
sur un port TCP. C'est aussi la plus dangereuse : cette API n'a, par défaut,
\textbf{aucune authentification}, et y accéder revient à être administrateur de
la machine. Les moteurs de recherche de services exposés en recensent des
dizaines de milliers, et ils sont compromis en quelques minutes.

\parade
Le transport SSH natif du client Docker. Aucun port supplémentaire n'est
ouvert : SSH était déjà là, et il apporte l'authentification que l'API n'a pas.

\danslecode
\begin{lstlisting}[language=Python]
def docker_command(worker: str = "local") -> list:
    """Prefixe de commande docker pour ce noeud."""
    host = worker_host(worker)
    if host is None:
        return ["docker"]
    return ["docker", "-H", f"ssh://{DOCKER_SSH_USER}@{host}"]
\end{lstlisting}

Le module le résume en tête : \emph{« Rien n'est exposé sur le réseau à part
SSH ; l'authentification se fait par clé. »}

\begin{notion}{Ce que fait réellement \cmd{docker -H ssh://}}
Le client Docker ouvre une session SSH et y exécute
\cmd{docker system dial-stdio} sur le nœud. L'API circule alors \emph{dans} le
tunnel SSH, authentifiée et chiffrée, sans qu'aucune socket réseau n'écoute.
Si l'on avait voulu sécuriser l'API TCP, il aurait fallu monter une autorité de
certification, émettre un certificat client par contrôleur et les renouveler :
beaucoup de pièces mobiles pour un résultat équivalent.
\end{notion}

\limite
Qui obtient la clé SSH du contrôleur obtient le démon Docker de tous les nœuds.
Cette clé est donc l'actif le plus sensible du parc --- d'où la mesure
suivante.

% -----------------------------------------------------------------------------
\mesure{La clé de pilotage}{mesure 25 · rôle \cmd{controller\_key}}

\menace
Une clé SSH réutilisée --- celle du développeur, par exemple --- étend sa portée
à tout le parc, et survit au départ de son propriétaire.

\parade
Une clé \textbf{dédiée au pilotage}, générée par Ansible, déposée sur les
nœuds pour un \textbf{compte de service sans mot de passe}.

\danslecode
Le rôle \cmd{controller\_key} crée le compte et la paire de clés ; le rôle
\cmd{worker} dépose la partie publique sur chaque nœud. Le compte
\cmd{insacloud} n'a pas de mot de passe --- il n'est donc joignable que par
clé, et uniquement depuis le contrôleur.

Sur le parc à deux nœuds en conteneurs, le même principe est rejoué par le
service \cmd{cles}, qui fabrique la paire une seule fois dans un volume
partagé. L'\fichier{entrypoint} du nœud installe la partie publique et, détail
qui a coûté une heure, déverrouille le compte :

\begin{lstlisting}[language=bash]
# « adduser -D » laisse le compte VERROUILLE (un « ! » dans /etc/shadow), et
# sshd refuse un compte verrouille MEME en authentification par cle. L'etoile
# dit « aucun mot de passe n'ouvre ce compte », ce qui est l'intention.
sed -i 's|^insacloud:!|insacloud:*|' /etc/shadow
\end{lstlisting}

\begin{piege}{Un compte sans mot de passe n'est pas un compte verrouillé}
Les deux états se ressemblent et n'ont rien à voir. \cmd{!} signifie
« verrouillé » : \cmd{sshd} refuse la connexion, y compris par clé.
\cmd{*} signifie « aucun mot de passe ne peut ouvrir ce compte » : la
connexion par clé fonctionne. Le symptôme était un
\cmd{Connection closed by remote host} muet côté client ; seul le journal du
serveur disait \cmd{User insacloud not allowed because account is locked}.
\end{piege}

\limite
La clé n'a pas de phrase de passe : elle doit être utilisable par un service,
sans intervention humaine. Sa protection repose donc entièrement sur les droits
du fichier et sur la sécurité du contrôleur.

% -----------------------------------------------------------------------------
\mesure{Le jeton des agents}{mesure 26 · \fichier{api.py}, \fichier{worker\_agent.py}}

\menace
L'inscription des nœuds est dynamique : un agent appelle
\cmd{POST /workers/register} et le contrôleur l'ajoute au parc. Sans contrôle,
n'importe qui déclare un nœud, reçoit des machines à héberger --- et donc les
mots de passe de leurs locataires.

\parade
Un jeton partagé, exigé dans un en-tête à l'inscription comme au battement de
cœur.

\danslecode
\begin{lstlisting}[language=Python]
"X-Agent-Token": TOKEN        # cote agent
\end{lstlisting}

Côté contrôleur, les routes \cmd{/workers/register} et \cmd{/workers/heartbeat}
refusent la requête sans jeton valide --- ce qu'un test vérifie
explicitement (\cmd{test\_api\_enregistrement\_refuse\_sans\_jeton}). Le jeton
est distribué par Ansible Vault, jamais inscrit dans le code.

L'agent, lui, refuse de démarrer sans jeton plutôt que de tenter une
inscription anonyme :

\begin{lstlisting}[language=Python]
if not TOKEN:
    log.error("INSACLOUD_AGENT_TOKEN n'est pas defini : arret.")
\end{lstlisting}

\limite
Un jeton partagé est le même pour tous les nœuds : il ne distingue pas les
agents entre eux, et le compromettre permet d'en déclarer un faux. Une
authentification mutuelle par certificat serait plus robuste ; elle demanderait
une autorité de certification que le projet n'a pas.

% -----------------------------------------------------------------------------
\mesure{Un nœud muet n'est jamais interprété}{mesure 27 · \fichier{faucheur.py}}

\menace
Le Faucheur réconcilie Docker et la base : il détruit les conteneurs inconnus
et marque arrêtées les machines sans conteneur. Si un nœud ne répond pas et
qu'il interprète ce silence comme « il n'y a plus rien là-bas », il déclare
mortes des machines qui tournent parfaitement --- et les détruira au retour du
nœud.

\parade
Un nœud injoignable est \textbf{ignoré}, pas interprété. Aucune décision n'est
prise sur son compte.

\danslecode
\begin{lstlisting}[language=Python]
if res.returncode != 0:
    log.warning("Noeud %s injoignable pour la reconciliation : %s",
                worker, res.stderr.strip()[:120])
    continue          # on ne touche a rien pour ce noeud
\end{lstlisting}

Un test dédié verrouille ce comportement :
\cmd{test\_noeud\_injoignable\_ne\_detruit\_rien} --- \emph{« Exigence de
résilience : un worker éteint ne doit pas faire perdre l'état en base. »}

\begin{retenir}
C'est une décision de sécurité autant que de fiabilité. La question posée au
système n'est pas « que s'est-il passé ? » mais « que puis-je affirmer ? ».
Un nœud silencieux ne permet d'affirmer rien du tout, et une opération
destructive exige une certitude.
\end{retenir}

\limite
La prudence a un coût : une machine réellement disparue reste marquée active
tant que son nœud ne répond pas. C'est le bon côté sur lequel se tromper.

% -----------------------------------------------------------------------------
\mesure{Le pare-feu, différencié par groupe}{mesure 28 · rôle \cmd{security\_hardening}}

\menace
Un nœud ouvre son port SSH pour le contrôleur, et des centaines de ports pour
les machines louées. Le contrôleur, lui, n'a besoin que de HTTPS. Appliquer la
même politique aux deux ouvre inutilement des portes.

\parade
Une politique par défaut en refus, puis des règles distinctes selon le groupe
Ansible auquel la machine appartient.

\danslecode
Le rôle \cmd{security\_hardening} applique, via \cmd{community.general.ufw} :
politique entrante par défaut, SSH avec \textbf{limitation anti-force-brute}
(\cmd{ufw limit}) et restriction d'origine, ports applicatifs, et la plage des
machines louées \emph{uniquement sur les workers}.

\begin{lstlisting}
- name: Autoriser SSH avec limitation anti-force-brute sur le port {{ ufw_ssh_port }}
  community.general.ufw:
    rule: limit
    port: "{{ ufw_ssh_port }}"

- name: Autoriser les plages de ports des machines louees
  community.general.ufw:
    ...
\end{lstlisting}

\begin{piege}{L'ordre des rôles n'est pas indifférent}
Le rôle passe \textbf{avant} le rôle \cmd{docker} : \cmd{ufw enable} recharge
\cmd{iptables} et effacerait les chaînes que Docker y installe. Inversé, le
pare-feu casserait le réseau des conteneurs --- panne difficile à diagnostiquer,
parce que les deux rôles réussissent.
\end{piege}

\limite
UFW filtre le trafic entrant vers l'hôte. Docker insère ses propres règles
\cmd{iptables} pour les ports publiés, qui peuvent contourner UFW selon la
configuration : la plage des machines louées est donc ouverte de fait dès que
Docker publie les ports.
